\section{Transport Layer, UDP, and Reliable Data Transfer}

\subsection{Overview of the Transport Layer}
The transport layer provides \textbf{process-to-process communication} on top of the network layer's host-to-host delivery. It enables multiple applications on the same host to communicate simultaneously using \textbf{ports} and \textbf{sockets}. :contentReference[oaicite:0]{index=0}

\textbf{Key responsibilities:}
\begin{itemize}
    \item Multiplexing: combining data from multiple applications
    \item Demultiplexing: delivering data to the correct process
    \item Providing logical communication between processes
\end{itemize}

\subsection{Ports and Sockets}
Each transport endpoint is identified using a \textbf{socket}, defined by a 5-tuple:
\[
(\text{protocol}, \text{local IP}, \text{local port}, \text{remote IP}, \text{remote port})
\]

\textbf{Port ranges:}
\begin{itemize}
    \item 0--1023: Well-known ports
    \item 1024--49151: Registered ports
    \item 49152--65535: Ephemeral ports
\end{itemize}

\subsection{User Datagram Protocol (UDP)}
UDP is a \textbf{connectionless} transport protocol that provides minimal services.

\textbf{Characteristics:}
\begin{itemize}
    \item No reliability guarantees
    \item No ordering or duplicate detection
    \item No congestion or flow control
    \item Low overhead (8-byte header)
\end{itemize}

\textbf{UDP Header Fields:}
\begin{itemize}
    \item Source Port
    \item Destination Port
    \item Length
    \item Checksum
\end{itemize}

\textbf{When to use UDP:}
\begin{itemize}
    \item Real-time applications (VoIP, gaming)
    \item Query-response protocols (DNS, NTP)
    \item Multicast/broadcast communication
    \item Applications implementing custom reliability (e.g., QUIC)
\end{itemize}

\subsection{UDP Checksum}
UDP uses a checksum for error detection over:
\begin{itemize}
    \item Pseudo-header (IP info)
    \item UDP header
    \item Data
\end{itemize}
It detects bit errors but does not guarantee delivery.

\subsection{Reliable Data Transfer (RDT)}
Reliable data transfer must handle:
\begin{itemize}
    \item Packet loss
    \item Packet corruption
\end{itemize}

\textbf{Key mechanisms:}
\begin{itemize}
    \item Checksums
    \item Acknowledgments (ACKs)
    \item Sequence numbers
    \item Timers
\end{itemize}

\subsection{Stop-and-Wait Protocol}
A simple ARQ protocol where the sender transmits one packet and waits for an ACK.

\textbf{Key properties:}
\begin{itemize}
    \item Only one packet in flight
    \item Uses sequence numbers (0/1) to detect duplicates
\end{itemize}

\textbf{Utilization:}
\[
U = \frac{t_{tx}}{RTT + t_{tx}}
\]

This leads to very poor efficiency when $RTT \gg t_{tx}$.

\subsection{Bandwidth-Delay Product (BDP)}
The bandwidth-delay product represents the number of bits in transit:
\[
BDP = B \times RTT
\]

To fully utilize a link, multiple packets must be in flight simultaneously.

\subsection{Pipelining}
Pipelining allows multiple unacknowledged packets to be sent, improving utilization.

\subsection{Go-Back-N (GBN)}
A sliding window protocol with the following properties:
\begin{itemize}
    \item Cumulative ACKs
    \item Single timer
    \item Receiver discards out-of-order packets
    \item Retransmits entire window on loss
\end{itemize}

\textbf{Constraint:}
\[
N \leq 2^k - 1
\]

\subsection{Selective Repeat (SR)}
An improved sliding window protocol:
\begin{itemize}
    \item Individual ACKs
    \item Multiple timers
    \item Receiver buffers out-of-order packets
    \item Retransmits only lost packets
\end{itemize}

\textbf{Constraint:}
\[
N \leq 2^{k-1}
\]

\subsection{Go-Back-N vs Selective Repeat}
\begin{itemize}
    \item GBN is simpler but inefficient under packet loss
    \item SR is more efficient but requires more complexity and buffering
\end{itemize}

\subsection{TCP Design Insights}
TCP combines ideas from both GBN and SR:
\begin{itemize}
    \item Uses cumulative ACKs (like GBN)
    \item Buffers out-of-order packets (like SR)
    \item Uses Selective Acknowledgment (SACK) for efficiency
\end{itemize}

\subsection{Looking Ahead}
A complete transport protocol must also handle:
\begin{itemize}
    \item Flow control
    \item Congestion control
    \item Connection setup and teardown
    \item Byte-stream abstraction
\end{itemize}